\begin{derivation}[从\eqnrefs{eq:qcd-covariant-derivative-dp68}到\eqnrefs{eq:qcd-field-strength-matrix-dp10}]
让 $[D_\mu,D_\nu]$ 作用在任意基本表示颜色场 $\chi$ 上。由于 $D_\mu$ 含微分算符，必须先完成对场的作用，再比较两种作用次序：

\begin{align*}
D_\mu D_\nu\chi
={}&\partial_\mu\partial_\nu\chi
+i g_s(\partial_\mu\mathscr C_\nu)\chi
+i g_s\mathscr C_\nu\partial_\mu\chi\\
&+i g_s\mathscr C_\mu\partial_\nu\chi
-g_s^2\mathscr C_\mu\mathscr C_\nu\chi,
\end{align*}
而交换 $\mu\leftrightarrow\nu$ 得
\begin{align*}
D_\nu D_\mu\chi
={}&\partial_\nu\partial_\mu\chi
+i g_s(\partial_\nu\mathscr C_\mu)\chi
+i g_s\mathscr C_\mu\partial_\nu\chi\\
&+i g_s\mathscr C_\nu\partial_\mu\chi
-g_s^2\mathscr C_\nu\mathscr C_\mu\chi.
\end{align*}
普通偏导彼此对易，且两组含 $\partial\chi$ 的交叉项逐项相消，所以
\begin{equation*}
[D_\mu,D_\nu]\chi
=\left\{i g_s(\partial_\mu\mathscr C_\nu-\partial_\nu\mathscr C_\mu)
-g_s^2[\mathscr C_\mu,\mathscr C_\nu]\right\}\chi.
\end{equation*}
提出 $i g_s$ 就得到\eqnrefs{eq:qcd-field-strength-matrix-dp10}。再写
\begin{equation*}
\mathscr C_\mu=\mathscr C_\mu^aT^a,
\qquad
[\mathscr C_\mu,\mathscr C_\nu]
=i f^{abc}\mathscr C_\mu^a\mathscr C_\nu^bT^c,
\end{equation*}
便得到分量形式中的 $-g_sf^{abc}\mathscr C_\mu^b\mathscr C_\nu^c$。
\end{derivation}

\begin{derivation}[从\eqnrefs{eq:qcd-field-strength-matrix-dp10}到\eqnrefs{eq:yang-mills-expanded-3-4-gluon-dp10}]
纯规范场部分的三次与四次相互作用可以从场强平方直接展开。定义

\begin{equation*}
\mathcal F_{\mu\nu}^{a(0)}\equiv
\partial_\mu\mathscr C_\nu^a-\partial_\nu\mathscr C_\mu^a,
\qquad
X_{\mu\nu}^a\equiv f^{abc}\mathscr C_\mu^b\mathscr C_\nu^c,
\end{equation*}
则 $\mathcal F_{\mu\nu}^a=\mathcal F_{\mu\nu}^{a(0)}-g_sX_{\mu\nu}^a$。平方为
\begin{equation*}
\mathcal F_{\mu\nu}^a\mathcal F^{a\mu\nu}
=\mathcal F_{\mu\nu}^{a(0)}\mathcal F^{a(0)\mu\nu}
-2g_s\mathcal F_{\mu\nu}^{a(0)}X^{a\mu\nu}
+g_s^2X_{\mu\nu}^aX^{a\mu\nu}.
\end{equation*}
乘上 $-\hbar c/4$ 后，依次得到自由项、三次项与四次项。三次项的系数还需要利用场强的反对称结构整理：
\begin{align*}
\mathcal F_{\mu\nu}^{a(0)}X^{a\mu\nu}
&=f^{abc}(\partial_\mu\mathscr C_\nu^a-\partial_\nu\mathscr C_\mu^a)
\mathscr C^{b\mu}\mathscr C^{c\nu}.
\end{align*}
第二项交换哑 Lorentz 指标 $\mu\leftrightarrow\nu$ 后变为
\begin{equation*}
-f^{abc}(\partial_\mu\mathscr C_\nu^a)
\mathscr C^{b\nu}\mathscr C^{c\mu}.
\end{equation*}
再交换哑颜色指标 $b\leftrightarrow c$，并利用 $f^{acb}=-f^{abc}$，它变成与第一项相同的
\begin{equation*}
f^{abc}(\partial_\mu\mathscr C_\nu^a)
\mathscr C^{b\mu}\mathscr C^{c\nu}.
\end{equation*}
因此两项相加产生因子 $2$，三次拉氏量可写成
\begin{equation*}
\mathcal L_{3g}
=\hbar c g_s f^{abc}(\partial_\mu\mathscr C_\nu^a)
\mathscr C^{b\mu}\mathscr C^{c\nu},
\end{equation*}
这就是\eqnrefs{eq:yang-mills-expanded-3-4-gluon-dp10} 的交叉项。四次项为
\begin{equation*}
\mathcal L_{4g}
=-\frac{\hbar c g_s^2}{4}
 f^{abc}f^{ade}
 \mathscr C_\mu^b\mathscr C_\nu^c
 \mathscr C^{d\mu}\mathscr C^{e\nu}.
\end{equation*}

\end{derivation}

